%% Architecture — \neven{} Paper (Restructured)

\section{System Architecture and Design}
\label{sec:architecture}

\neven{} is organized as a four-layer architecture (Figure~\ref{fig:architecture}) that separates concerns between the Excel interface, core services, language runtimes, and the security subsystem. This section describes each layer and the inter-process communication protocol that connects them.

\subsection{Architectural Overview}

\begin{figure}[htbp]
    \centering
	\includegraphics[width=0.99\textwidth]{sections/architecture_diagram.tex}
    \caption{\neven{} four-layer architecture. The XLL add-in (Layer~1) communicates with language runtimes via Named Pipes carrying Protocol Buffer messages. Each runtime executes in an isolated process, preventing fault propagation to the host application. WebView2 provides interactive visualization on a dedicated STA thread.}
    \label{fig:architecture}
\end{figure}

The design follows three guiding principles:

\begin{enumerate}
    \item \textbf{Fault isolation:} Language runtimes execute in separate OS processes. A segmentation fault in R, Julia, or Python terminates only the child process; the Excel host remains operational and can reconnect automatically.
    \item \textbf{Protocol neutrality:} All data exchange uses Protocol Buffers \citep{Protobuf2023}, providing language-agnostic serialization with backward-compatible schema evolution.
    \item \textbf{Security by default:} User-submitted code passes through a sandbox verifier before execution, blocking dangerous operations at the string-analysis level.
\end{enumerate}

\subsection{System Initialization}
\label{subsec:initialization}

The entry point is \texttt{NEVEN64.xll}, a 64-bit XLL add-in compiled as a dynamic-link library. Upon loading, Excel invokes \texttt{xlAutoOpen()}, which triggers the initialization sequence detailed in Algorithm~\ref{alg:init}. The XLL registers approximately 200 functions with the Excel function wizard via the \texttt{xlfRegister} API call.

\begin{algorithm}[htbp]
\SetAlgoLined
\KwInput{XLL loaded by Excel}
\KwOutput{System ready with all functions registered}
\BlankLine
\texttt{ZombieProcessKiller()}\Comment*[r]{Kill orphaned child processes}
\texttt{ConfigService::Initialize}(``neven-config.json'')\;
\texttt{DiscoveryService::FindLanguages()}\;
\BlankLine
\ForEach{language $L \in \{R, Julia, Python\}$}{
    \If{$L$ is enabled and discovered}{
        $proc \leftarrow$ \texttt{CreateProcess}($L$.executable)\;
        $pipe \leftarrow$ \texttt{ConnectNamedPipe}(``RJ2XCL2-PIPE-$L$-'' + PID)\;
        \texttt{SendStartupScript}($pipe$, $L$.startup\_file)\;
        \texttt{LoadUserFunctions}($pipe$, $L$.functions\_dir)\;
    }
}
\BlankLine
$functions \leftarrow$ \texttt{MapFunctions()}\Comment*[r]{Query all runtimes}
\ForEach{$f \in functions$}{
    \texttt{xlfRegister}($f$.name, $f$.params, $f$.description)\;
}
\texttt{ViewerManager::Initialize()}\Comment*[r]{WebView2 STA thread}
\texttt{PlutoManager::Initialize()}\;
\caption{System Initialization Sequence (\texttt{xlAutoOpen})}
\label{alg:init}
\end{algorithm}

During the initialization sequence shown in Algorithm~\ref{alg:init}, user-created functions are automatically discovered and registered alongside the built-in library. Users extend the system by placing \texttt{.R}, \texttt{.jl}, or \texttt{.py} files in the designated functions directory (within the user's \texttt{Documents\textbackslash NEVEN\textbackslash functions\textbackslash} folder). These files are loaded into their respective runtimes during the \texttt{LoadUserFunctions} step, and every function defined therein is registered as a native Excel formula. The user simply restarts Excel to make new functions available---no compilation, no configuration editing, no administrative privileges required.

The XLL implements two function dispatch mechanisms:
\begin{itemize}
    \item \textbf{Registered functions} (\texttt{=R.MR\_Lineal(...)}, \texttt{=J.Algebra(...)}): Dispatched via \texttt{RJ\_FunctionCall}, which serializes arguments to Protocol Buffers and routes to the appropriate language service.
    \item \textbf{Generic execution} (\texttt{=NEVEN.r("code")}, \texttt{=NEVEN.j("code")}): Dispatched via \texttt{RJ\_Exec\_Generic}, which first validates code through the sandbox verifier before forwarding to the language runtime.
\end{itemize}

A COM-based Ribbon add-in (\texttt{NEVENRibbon.dll}) provides a native Excel tab with 13 buttons organized in five groups, offering direct access to language consoles, visualization tools, Pluto.jl notebooks, Quarto rendering, and system configuration.

\subsection{Core Services}

\paragraph{Configuration Service.}
The \texttt{ConfigService} reads \texttt{neven-config.json} at startup, providing typed accessors for all operational parameters. Configuration validation rejects paths containing traversal sequences (\texttt{..}), command injection characters (\texttt{|}, \texttt{\&}, \texttt{;}), and enforces numeric ranges for timeouts and retry counts.

\paragraph{Language Manager.}
The \texttt{LanguageManager} orchestrates connections to all language runtimes. For each configured language, it launches the child process via \texttt{CreateProcessA}, establishes a Named Pipe connection with the naming pattern \texttt{RJ2XCL2-PIPE-\{LANG\}-\{PID\}}, sends the startup script synchronously, loads user function files from the configured directory, and queries the runtime for its registered function list. Health monitoring runs on a 5-second timer, checking process exit codes via \texttt{GetExitCodeProcess}. If a runtime becomes unavailable, the system marks it as disconnected and returns descriptive error messages rather than opaque \texttt{\#VALUE!} errors.

\paragraph{Discovery Service.}
The \texttt{DiscoveryService} locates R, Julia, and Python installations by querying the Windows Registry (\texttt{HKLM\textbackslash Software\textbackslash R-core}), scanning standard installation paths, and checking the system \texttt{PATH}. Version constraints (minimum and maximum major/minor versions) are enforced from configuration, ensuring compatibility with tested runtime versions.

\subsection{Inter-Process Communication Protocol}
\label{subsec:ipc}

Communication between the XLL and language runtimes uses Windows Named Pipes with Protocol Buffers serialization. The message exchange protocol is formalized in Algorithm~\ref{alg:ipc}, which describes the complete lifecycle of a function call from Excel cell to language runtime and back.

\begin{algorithm}[htbp]
\SetAlgoLined
\KwInput{Function name $f$, arguments $\{a_1, \ldots, a_n\}$ from Excel cells}
\KwOutput{Result matrix $R$ for Excel display}
\BlankLine
\Comment{Phase 1: Serialization}
$msg \leftarrow$ \texttt{Variable.new()}\;
$msg.name \leftarrow f$\;
\ForEach{argument $a_i$}{
    \eIf{$a_i$ is scalar}{
        $msg.value[i] \leftarrow$ \texttt{serialize\_scalar}($a_i$)\;
    }{
        $msg.value[i] \leftarrow$ \texttt{serialize\_range}($a_i$)\Comment*[r]{Row-major matrix}
    }
}
$payload \leftarrow$ \texttt{msg.SerializeToString()}\;
\BlankLine
\Comment{Phase 2: Pipe Communication}
$pipe \leftarrow$ \texttt{GetPipe}(language)\;
\texttt{WriteFile}($pipe$, \texttt{length}($payload$))\Comment*[r]{4-byte length prefix}
\texttt{WriteFile}($pipe$, $payload$)\;
\BlankLine
\Comment{Phase 3: Response}
$resp\_len \leftarrow$ \texttt{ReadFile}($pipe$, 4)\;
$resp\_data \leftarrow$ \texttt{ReadFile}($pipe$, $resp\_len$)\;
$result \leftarrow$ \texttt{Variable.ParseFromString}($resp\_data$)\;
\BlankLine
\Comment{Phase 4: Deserialization to XLOPER12}
\eIf{$result$ is error}{
    \Return \texttt{xloper\_string}($result.error\_message$)\;
}{
    \Return \texttt{to\_xloper12}($result$)\Comment*[r]{Matrix or scalar}
}
\caption{IPC Protocol for Function Dispatch}
\label{alg:ipc}
\end{algorithm}

The Protocol Buffers schema (\texttt{variable.proto}) defines a \texttt{Variable} message type that supports scalars (real, integer, string, boolean, nil), arrays, and matrices with explicit dimension metadata. This schema has remained stable since version~1.0, demonstrating the forward-compatibility benefits of Protocol Buffers for long-lived IPC interfaces.

\paragraph{Performance Characteristics.}
The Named Pipe transport operates in byte mode with a buffer size of 65,536 bytes. For typical statistical function calls (arguments under 10~KB), the serialization--transmission--deserialization round trip adds less than 5~ms of overhead beyond the computation time in the target language. Larger payloads (matrices exceeding 100,000 cells) are transmitted in chunked writes to avoid pipe buffer exhaustion.

\paragraph{Julia Cold Start Elimination.}
Julia's just-in-time compilation introduces significant latency on first function calls. \neven{} addresses this through a precompiled system image (\texttt{neven\_julia.dll}, approximately 415~MB) generated via \texttt{PackageCompiler.jl} \citep{PackageCompiler2023}. This image includes all \neven{} Julia modules pre-compiled to native code, reducing first-call latency from minutes to under one second.

\subsection{Security Sandbox}
\label{subsec:sandbox}

The \texttt{SandboxVerifier} implements a pattern-based security layer that validates all user-submitted code before execution. Table~\ref{tab:sandbox} summarizes the blocked categories.

\begin{table}[htbp]
\centering
\small
\caption{Security sandbox: blocked operation categories per language. Each category includes multiple specific patterns (30+ total per language).}
\label{tab:sandbox}
\resizebox{\textwidth}{!}{%
\begin{tabular}{@{}lp{3cm}p{3cm}p{3cm}@{}}
\toprule
\textbf{Category} & \textbf{R} & \textbf{Julia} & \textbf{Python} \\
\midrule
Shell execution & \texttt{system()}, \texttt{shell()}, \texttt{pipe()} & \texttt{run()}, \texttt{pipeline()}, backticks & \texttt{os.system()}, \texttt{subprocess.*} \\
\addlinespace
File system & \texttt{file.remove()}, \texttt{unlink()} & \texttt{rm()}, \texttt{mv()}, \texttt{cp()} & \texttt{os.remove()}, \texttt{shutil.*} \\
\addlinespace
Network & \texttt{download.file()}, \texttt{url()} & \texttt{download()} & \texttt{urllib.*}, \texttt{requests.*} \\
\addlinespace
Dynamic code & \texttt{eval(parse())}, \texttt{do.call()} & \texttt{eval()}, \texttt{Meta.parse()} & \texttt{exec()}, \texttt{eval()}, \texttt{compile()} \\
\addlinespace
Native code & \texttt{.Call()}, \texttt{dyn.load()} & \texttt{ccall()}, \texttt{unsafe\_*} & \texttt{ctypes.*}, \texttt{cffi.*} \\
\bottomrule
\end{tabular}%
}
\end{table}

The sandbox implements five bypass-prevention mechanisms: (1)~\emph{whitespace normalization}, which strips extraneous whitespace before pattern matching to prevent evasion via constructions like \texttt{sys\allowbreak tem()}; (2)~\emph{concatenation detection}, which identifies string-building patterns that assemble blocked function names at runtime; (3)~\emph{case-insensitive matching} to prevent \texttt{SYSTEM()} or \texttt{System()} evasion; (4)~\emph{context-aware detection}, which identifies dangerous function calls only in specific argument contexts (e.g., \texttt{readLines(pipe(...))} is blocked while \texttt{readLines("file.txt")} is allowed); and (5)~\emph{unified enforcement}, ensuring identical security checks regardless of execution source (Excel cell, REPL console, or AutoLoader script).

Additionally, a centralized \texttt{InputSanitizer} validates all file paths and command-line arguments against a character allowlist before any \texttt{CreateProcess} call, and a \texttt{MessageValidator} validates IPC message frames before Protobuf deserialization, providing defense-in-depth across the input, execution, and communication layers.

Registered functions---those loaded from \texttt{.R}, \texttt{.jl}, and \texttt{.py} files at startup---bypass the sandbox entirely, as they are considered trusted code authored by the system administrator or the user themselves.

\subsection{Visualization Subsystem}

\neven{} integrates Microsoft WebView2 \citep{WebView2_2023} (the Chromium-based rendering engine from Microsoft Edge) as an embedded visualization viewer. The \texttt{ViewerManager} singleton operates on a dedicated Single-Threaded Apartment (STA) COM thread, managing up to 8 concurrent viewer windows.

The visualization pipeline supports multiple rendering libraries:
\begin{itemize}
    \item \textbf{Plotly.js} \citep{Plotly2023}: Interactive scatter, bar, line, box, histogram, and heatmap charts with axis selectors.
    \item \textbf{D3.js v7} \citep{Bostock2011}: Treemap, Sankey diagram, Sunburst, and force-directed graph visualizations.
    \item \textbf{Leaflet.js} \citep{Leaflet2023}: Interactive maps with markers, heatmaps, and proportional circles.
    \item \textbf{rpivotTable} \citep{rpivotTable2023}: Drag-and-drop pivot tables with integrated charting.
\end{itemize}

A \texttt{PostMessageBridge} enables bidirectional communication between the JavaScript layer in WebView2 and the C++ host, allowing visualizations to request data or trigger Excel operations.

\subsection{Interactive REPL Console}
\label{subsec:repl}

\neven{} provides an interactive REPL (Read-Eval-Print Loop) console for exploratory coding, accessible via \texttt{=NEVEN.Console()} or the Ribbon button. The console is implemented as a self-contained HTML page rendered in WebView2, managed by a \texttt{REPLManager} singleton and a \texttt{REPLBridge} that handles code dispatch, sandbox enforcement, and result delivery.

This architecture replaces the legacy Electron-based console (inherited from BERT) that carried over 50 known CVEs from Electron~1.8.2, required Node.js and npm as dependencies, and operated as a separate 200~MB child process. The WebView2-based replacement introduces zero external dependencies (the WebView2 runtime is pre-installed on Windows~10/11), runs in-process on a dedicated STA thread, and applies the same \texttt{SandboxVerifier} checks as Excel cell execution---ensuring that the REPL cannot bypass security restrictions that apply to formulas.
